\subsection{Strategies for AI Safety and Risk Mitigation}
\label{sec:6.3.3}
Four practices do most of the work, and each is a way of not committing to a specification you cannot check.

Learn the objective instead of fixing it. Cooperative inverse reinforcement learning treats the human reward as unknown and to be inferred through interaction rather than handed over once; section 4.2.2 introduces it and section 5.6.3 works through what it costs. A system that keeps revising its model of what is wanted is safer than one built against a specification written before anyone had used it.

Quantify uncertainty. Bayesian methods and ensembles give a system a way to represent what it does not know, and a system that can recognize an input unlike anything in its training data is in a different position from one that classifies everything with the same unearned confidence.

Constrain the search, not the result. Constrained policy optimization keeps a reinforcement-learning agent inside its safety limits while it is looking for a better policy, rather than checking for violations afterward \autocite{achiam2017constrained} — the difference between a robot arm that learns to handle fragile objects without breaking one and a robot arm that learns by breaking several.

Keep the oversight live. Monitoring, anomaly detection, and a reviewer with the standing to override are deployment properties rather than training properties, and section 8.6.4 is the instrument for telling whether that standing is real or nominal.

Explanation belongs here too, as infrastructure rather than courtesy: a system that can say why it decided is a system whose failures are visible. Section 6.3.4 covers how the explanation gets produced and how much of it to believe.
